\appendix
\section{Compact formula atlas}
\label{app:atlas}

This appendix collects the formulas in one place.  Every quotient carries the
apartness conditions explained in the body, and every square root is the
positive branch justified by properness.  Put
\[
 p=\frac{a+b+c}{2},\qquad P=\frac{A+B+C}{2},\qquad
 E=A+B+C-\pi.
\]

\subsection{Primary laws}

\begin{align*}
\cos a&=\cos b\cos c+\sin b\sin c\cos A,\\
\cos b&=\cos c\cos a+\sin c\sin a\cos B,\\
\cos c&=\cos a\cos b+\sin a\sin b\cos C;\\[0.3em]
\cos A&=-\cos B\cos C+\sin B\sin C\cos a,\\
\cos B&=-\cos C\cos A+\sin C\sin A\cos b,\\
\cos C&=-\cos A\cos B+\sin A\sin B\cos c;\\[0.3em]
\frac{\sin A}{\sin a}&=
\frac{\sin B}{\sin b}=
\frac{\sin C}{\sin c}.
\end{align*}

Rearrangements:
\begin{align*}
\cos A&=\frac{\cos a-\cos b\cos c}{\sin b\sin c},&
\cos a&=\frac{\cos A+\cos B\cos C}{\sin B\sin C},
\end{align*}
plus cyclic versions.

\subsection{Polar triangle}

\begin{align*}
a^*&=\pi-A,&b^*&=\pi-B,&c^*&=\pi-C,\\
A^*&=\pi-a,&B^*&=\pi-b,&C^*&=\pi-c,
\end{align*}
and polarity is involutive.

\subsection{Half-angle and half-side families}

For \((A,a)\), with cyclic companions:
\begin{align*}
\sin\frac A2&=\sqrt{\frac{\sin(p-b)\sin(p-c)}{\sin b\sin c}},\\
\cos\frac A2&=\sqrt{\frac{\sin p\sin(p-a)}{\sin b\sin c}},\\
\tan\frac A2&=\sqrt{\frac{\sin(p-b)\sin(p-c)}{\sin p\sin(p-a)}};\\[0.4em]
\sin\frac a2&=\sqrt{-\frac{\cos P\cos(P-A)}{\sin B\sin C}},\\
\cos\frac a2&=\sqrt{\frac{\cos(P-B)\cos(P-C)}{\sin B\sin C}},\\
\tan\frac a2&=\sqrt{-\frac{\cos P\cos(P-A)}{\cos(P-B)\cos(P-C)}}.
\end{align*}

Common factors:
\begin{align*}
k&=[\sin p\sin(p-a)\sin(p-b)\sin(p-c)]^{-1/2},\\
\tan\frac A2&=k\sin(p-b)\sin(p-c),\\
K&=\sqrt{-\frac{\cos P}{\cos(P-A)\cos(P-B)\cos(P-C)}},\\
\tan\frac a2&=K\cos(P-A),
\end{align*}
plus cyclic companions.

Direct vector forms, where
\(x=\cos a,y=\cos b,z=\cos c\) and
\(\tau=\detm[u\ v\ w]\):
\begin{align*}
\tan\frac A2&=\frac{\tau}{\sin b\sin c+\cos a-\cos b\cos c},\\
\tan\frac a2&=\frac{\sin a}{1+\cos a}.
\end{align*}

With \(t_a=\tan(a/2)\) and similarly for all six parts:
\begin{align*}
t_A^2&=
 \frac{(t_a-t_b+t_c+t_at_bt_c)(t_a+t_b-t_c+t_at_bt_c)}
 {(t_a+t_b+t_c-t_at_bt_c)(-t_a+t_b+t_c+t_at_bt_c)},\\
t_a^2&=
 \frac{(t_Bt_C+t_Ct_A+t_At_B-1)
       (-t_Bt_C+t_Ct_A+t_At_B+1)}
 {(t_Bt_C-t_Ct_A+t_At_B+1)
       (t_Bt_C+t_Ct_A-t_At_B+1)}.
\end{align*}

\subsection{Haversines}

\begin{align*}
\hav a&=\hav(b-c)+\sin b\sin c\,\hav A,\\
\hav b&=\hav(c-a)+\sin c\sin a\,\hav B,\\
\hav c&=\hav(a-b)+\sin a\sin b\,\hav C;\\
\hav A&=\hav(\pi-B-C)+\sin B\sin C\,\hav a,
\end{align*}
with the two cyclic duals and guarded inverse forms.

\subsection{Delambre--Gauss analogies}

\begin{align*}
\frac{\sin((A+B)/2)}{\cos(C/2)}&=
 \frac{\cos((a-b)/2)}{\cos(c/2)},\\
\frac{\sin((A-B)/2)}{\cos(C/2)}&=
 \frac{\sin((a-b)/2)}{\sin(c/2)},\\
\frac{\cos((A+B)/2)}{\sin(C/2)}&=
 \frac{\cos((a+b)/2)}{\cos(c/2)},\\
\frac{\cos((A-B)/2)}{\sin(C/2)}&=
 \frac{\sin((a+b)/2)}{\sin(c/2)},
\end{align*}
plus eight cyclic companions.  The cross-multiplied forms are the primary
constructive statements.

\subsection{Napier analogies and tangent law}

\begin{align*}
\tan\frac{A+B}{2}&=
 \frac{\cos((a-b)/2)}{\cos((a+b)/2)}\cot\frac C2,\\
\tan\frac{A-B}{2}&=
 \frac{\sin((a-b)/2)}{\sin((a+b)/2)}\cot\frac C2,\\
\tan\frac{a+b}{2}&=
 \frac{\cos((A-B)/2)}{\cos((A+B)/2)}\tan\frac c2,\\
\tan\frac{a-b}{2}&=
 \frac{\sin((A-B)/2)}{\sin((A+B)/2)}\tan\frac c2,
\end{align*}
plus cyclic companions, and
\[
\frac{\tan((A-B)/2)}{\tan((A+B)/2)}
 =\frac{\tan((a-b)/2)}{\tan((a+b)/2)}
\]
plus its two cyclic companions.

\subsection{Cotangent and five-part formulas}

\begin{align*}
\cos b\cos C&=\cot a\sin b-\cot A\sin C,\\
\cos b\cos A&=\cot c\sin b-\cot C\sin A,\\
\cos c\cos A&=\cot b\sin c-\cot B\sin A,\\
\cos c\cos B&=\cot a\sin c-\cot A\sin B,\\
\cos a\cos B&=\cot c\sin a-\cot C\sin B,\\
\cos a\cos C&=\cot b\sin a-\cot B\sin C.
\end{align*}
Representative five-part forms:
\begin{align*}
\cos a\sin c&=\sin a\cos c\cos B+\sin b\cos A,\\
\sin a\cos B&=\cos b\sin c-\sin b\cos c\cos A,
\end{align*}
with cyclic and polar companions.

Cagnoli:
\[
\sin b\sin c+\cos b\cos c\cos A
 =\sin B\sin C-\cos B\cos C\cos a,
\]
plus two cyclic companions.

\subsection{Right triangle, \(C=\pi/2\)}

\begin{align*}
\cos c&=\cos a\cos b,&
\sin a&=\sin A\sin c,&
\sin b&=\sin B\sin c,\\
\tan a&=\tan A\sin b,&
\tan b&=\tan B\sin a,&
\tan b&=\cos A\tan c,\\
\tan a&=\cos B\tan c,&
\cos A&=\sin B\cos a,&
\cos B&=\sin A\cos b,\\
\cos c&=\cot A\cot B.&&
\end{align*}

\subsection{Quadrantal triangle, \(c=\pi/2\)}

\begin{align*}
\cos C&=-\cos A\cos B,&
\sin A&=\sin a\sin C,&
\sin B&=\sin b\sin C,\\
\tan A&=\tan a\sin B,&
\tan B&=\tan b\sin A,&
\tan B&=-\cos a\tan C,\\
\tan A&=-\cos b\tan C,&
\cos a&=\sin b\cos A,&
\cos b&=\sin a\cos B,\\
\cos C&=-\cot a\cot b.&&
\end{align*}

\subsection{Differential, area, and solid-angle formulas}

\begin{align*}
da&=\cos C\,db+\cos B\,dc+\sin b\sin C\,dA,\\
db&=\cos A\,dc+\cos C\,da+\sin c\sin A\,dB,\\
dc&=\cos B\,da+\cos A\,db+\sin a\sin B\,dC.\\[0.3em]
\Area(ABC)&=E=A+B+C-\pi,\\
\Area(P_n)&=\sum_{j=1}^n\alpha_j-(n-2)\pi,\\
E&=2\atanTwo\bigl(\tau,1+\cos a+\cos b+\cos c\bigr),\\
\tan^2\frac E4&=
 \tan\frac p2\tan\frac{p-a}{2}\tan\frac{p-b}{2}\tan\frac{p-c}{2}.
\end{align*}

\subsection{Tetrahedral formulas}

For \(\{i,j,k,l\}=\{1,2,3,4\}\):
\begin{align*}
\Delta_i^2={}&\Delta_j^2+\Delta_k^2+\Delta_l^2\\
&-2\bigl(\Delta_j\Delta_k\cos\theta_{il}
+\Delta_j\Delta_l\cos\theta_{ik}
+\Delta_k\Delta_l\cos\theta_{ij}\bigr),\\
\Delta_i={}&\Delta_j\cos\theta_{kl}
+\Delta_k\cos\theta_{jl}
+\Delta_l\cos\theta_{jk}.
\end{align*}
For vertices \(O,A,B,C\):
\[
\sin\angle OAB\sin\angle OBC\sin\angle OCA
=\sin\angle OAC\sin\angle OCB\sin\angle OBA.
\]
